% The unit-cube orientation trap from section 8.8.
% Shows the cube; on the bottom face the induced normal points inward (red);
% on the top face it points outward (blue, correct).
\begin{tikzpicture}[>=Stealth, scale=1.4]
  % cube vertices
  \coordinate (A) at (0,0);
  \coordinate (B) at (2,0);
  \coordinate (C) at (2.7,0.7);
  \coordinate (D) at (0.7,0.7);
  \coordinate (E) at (0,2);
  \coordinate (F) at (2,2);
  \coordinate (G) at (2.7,2.7);
  \coordinate (H) at (0.7,2.7);

  % back faces (dashed)
  \draw[dashed, enggray] (A)--(D)--(C);
  \draw[dashed, enggray] (D)--(H);

  % front + top + right faces
  \draw[thick] (A)--(B)--(F)--(E)--cycle;
  \draw[thick] (B)--(C)--(G)--(F);
  \draw[thick] (E)--(H)--(G);

  % top-face outward normal (correct, blue, up)
  \draw[->, very thick, engblue]
    ($(E)!0.5!(G)$) -- ++(0,0.9) node[above, font=\footnotesize]{$+\unitvec{z}$ outward (top, correct)};

  % bottom-face induced normal (wrong, red, also up but pointing into cube)
  \draw[->, very thick, engred]
    ($(A)!0.5!(C)$) -- ++(0,0.9) node[above right, font=\footnotesize, align=left]
    {$+\unitvec{z}$ induced (bottom)\\points \emph{into} the cube};

  % labels
  \node[font=\footnotesize] at ($(A)!0.5!(B)+(0,-0.18)$) {bottom face $z=0$};
  \node[font=\footnotesize] at ($(H)!0.5!(G)+(0,0.18)$) {top face $z=1$};
\end{tikzpicture}
